\section{Decision Making with Videos}\label{sec:method}
\input{figText/method_overview.tex}



Next we describe the proposed approach \model in detail, which is a concrete instantiation of the diffusion \proc. \model incorporates %
the two main components discussed in Section~\ref{sec:background}, %
as shown in~\fig{fig:method_overview}: {\bf (i)} a diffusion model for the universal video-based planner $\rho(\cdot|x_0, c)$, which synthesizes videos conditioned on the first frame and task descriptions; and {\bf (ii)} a task-specific action generator $\pi(\cdot|\{x_h\}_{h=0}^H, c)$, which infers actions sequences from generated videos through inverse dynamics modeling.



\subsection{Universal Video-Based Planner}

Encouraged by the recent success of text-to-video models \citep{ho2022imagen}, we seek to construct a video diffusion module as the trajectory planner, which can faithfully synthesize future image frames given an initial frame and textual task description. 
However, the desired planner departs from the typical setting in text-to-video models \citep{villegas2022phenaki, ho2022imagen} which normally generate unconstrained videos given a text description. Planning through video generation is more challenging as it requires models to both be able to generate constrained videos that start at a specified image, and then complete the target task. 
Moreover, to ensure valid action inference across synthesized frames in a video, the video prediction module needs to be able to track the underlying environment state across synthesized video frames. %

\myparagraph{Conditional Video Synthesis.} To generate a valid and executable plan, a text-to-video model must synthesize a constrained video plan starting at an initial image that depicts the initial configuration of the agent and environment. %
One approach to solve this problem is to modify the underlying \emph{test-time} sampling procedure of an unconditional model, by fixing the first frame of the generated video plan to always begin at the observed image, as done in ~\citep{janner2022planning}. However, we found that this performed poorly and led to subsequent frames in the video plan to deviate significantly from the original observed image. Instead, we found it more effective to explicitly train a constrained video synthesis model by providing the first frame of each video as explicit conditioning context during training. 

\myparagraph{Trajectory Consistency through Tiling.} 
Existing text-to-video models typically generate videos where the underlying environment state changes significantly during the temporal duration \citep{ho2022imagen}. To construct an accurate trajectory planner, it is important that the environment remain consistent across all time points. To enforce environment consistency in conditional video synthesis, we provide, as additional context, the observed image when denoising each frame in the synthesized video. 
In particular, we re-purpose a temporal super-resolution video diffusion architecture, and provide as context the conditioned visual observation tiled across time, as the opposed to a low temporal-resolution video for denoising at each timestep. In this model, we directly concatenate each intermediate noisy frame with the conditioned observed image across sampling steps, which serves as a strong signal to maintain the underlying environment state across time.



\looseness=-1
\myparagraph{Hierarchical Planning.} When constructing plans in high dimensional environments with long time horizons, directly generating a set of actions to reach a goal state quickly becomes intractable due to the exponential blow-up of the underlying search space. Planning methods often circumvent this issue by leveraging a natural hierarchy in planning. Specifically, planning methods first construct coarse plans operating on low dimensional states and actions, which may then be refined into plans in the underlying state and action spaces.
Similar to planning, our conditional video generation procedure likewise exhibits a natural temporal hierarchy. We first generate videos at a coarse level by sparsely sampled videos (``abstractions'') of the desired behavior along the time axis. Then we refine the videos to represent valid behavior in the environment by super-resolving videos across time. Meanwhile, coarse-to-fine super-resolution further improves consistency via interpolation between frames. 

\myparagraph{Flexible Behavioral Modulation.} 
When planning a sequence of actions to a given sub-goal, one can readily incorporate external constraints to modulate the generated plan. Such test-time adaptability can be implemented by composing a prior $h(\tau)$ during plan generation to specify desired constraints across the synthesized action trajectory \citep{janner2022planning},
which is also compatible with \model. 
In particular, the prior $h(\tau)$ can be specified using a learned classifier on images to optimize a particular task, or as a Dirac delta on a particular image to guide a plan towards a particular set of states.
To train the text-conditioned video generation model, we utilize the video diffusion algorithm in \citep{ho2022imagen}, where pretrained language features from T5 \citep{raffel2020exploring} are encoded. Please see Appendix~\ref{app:details} for the underlying architecture and training details.

\subsection{Task Specific Action Adaptation}

Given a set of synthesized %
videos,
we may train a small task-specific inverse-dynamics model to translate frames into a set of actions as described below.

\myparagraph{Inverse Dynamics.} We train a small model to estimate actions given input images. The training of the inverse dynamics is independent from the planner and can be done on a separate, smaller and potentially suboptimal dataset generated by a simulator.

\myparagraph{Action Execution.} 
Finally, we generate an action sequence given $x_0$ and $c$ by synthesizing $H$ image frames and applying the learned inverse-dynamics model to predict the corresponding $H$ actions.
Inferred actions can then be executed via \emph{closed-loop} control,
where we generate $H$ new actions after each step of action execution (i.e., model predictive control), 
or via \emph{open-loop} control, where we sequentially execute each action from the intially inferred action sequence. 
For computational efficiency, we use an open-loop controller in all our experiments in this paper.

